\uchapter{Abstract}

Decision support information systems are built on data warehouses, and two pressures now
strain them: growing source heterogeneity and rising analytical ambition. The literature
studies the resulting concerns separately --- \textit{structural scalability}, the capacity
to absorb schema change and new sources without disruptive re-engineering, and
\textit{analytical scalability}, the capacity to support the full analytical maturity
spectrum, from business intelligence to prediction, natively in the pipeline.

This thesis is a state-of-the-art study. It surveys and compares the principal architectures
on each axis against five criteria drawn from the literature, and proposes no architecture of
its own. On the structural axis it compares the Inmon Corporate Information Factory, the Kimball
dimensional model, Data Vault~2.0, and Anchor modeling, and finds that Data Vault~2.0 best
absorbs change because it makes schema evolution additive rather than destructive. On the
analytical axis it compares the Lambda and Kappa designs, the data lake, the Lakehouse, and
the Delta architecture, and finds that the Lakehouse, implemented through the Delta
architecture, best supports native, full-spectrum analytics. It then reviews the existing
combination of the two --- Data Vault~2.0 on the Lakehouse --- and finds it established in
industry practice but thinly evaluated in peer-reviewed work. Every conclusion traces to the
surveyed literature.

\bigskip
\vspace{0.5cm}

\noindent\textbf{Keywords:} Data Vault~2.0, Lakehouse, Decision support system,
Structural scalability, Analytical scalability, State of the art.
